%% =========================================================
%% Chapter: Schur labels of tensor powers
%% Group-algebra labels, permutation representations, and the Schur-label
%% estimates of the area-law paper (Lemma 6.1, lem:schur).
%% =========================================================

\chapter{Schur Labels of Tensor Powers}\label{ch:schur_labels}

This chapter develops the label calculus for permutations of the copies of a
finite-dimensional tensor power that is used for symmetric replicas in the
two-dimensional area-law argument~\cite{OpenAI2026AreaLaw}. The labels are the
irreducible representations of the symmetric group $S_k$; we read them off
from an Artin--Wedderburn decomposition of the complex group algebra, which
treats every finite group in the same way. The identification of the labels of
$S_k$ with partitions of $k$ enters only in the dimension and branching formulas.

\section{Labels of a finite group}

\begin{definition}[Irreducible labels and central idempotents]
    \label{def:irrep_labels}
    \lean{IrrepLabel, IrrepLabel.dim, IrrepLabel.wedderburnEquiv,
        IrrepLabel.centralIdem, IrrepLabel.matrixUnit}
    \leanok
    Let $G$ be a finite group. By Maschke's theorem and the Artin--Wedderburn
    theorem over $\C$ there is an algebra isomorphism
    \[
        \C[G]\;\cong\;\prod_{\lambda}M_{d_\lambda}(\C),
    \]
    and we fix one. Its factors are the labels $\lambda$ of $G$, and
    $d_\lambda\ge1$ is the dimension of the irreducible representation with
    label $\lambda$. The central idempotent $e_\lambda\in\C[G]$ and the matrix
    units $e^\lambda_{ij}$ are the preimages of the identity and of the matrix
    units of the factor $M_{d_\lambda}(\C)$.
\end{definition}

\begin{lemma}[Central idempotents]
    \label{lem:irrep_central_idempotents}
    \lean{IrrepLabel.centralIdem_mul_centralIdem, IrrepLabel.sum_centralIdem,
        IrrepLabel.centralIdem_mul_comm, IrrepLabel.centralIdem_ne_zero,
        IrrepLabel.invStar_centralIdem}
    \leanok
    \uses{def:irrep_labels}
    The central idempotents are nonzero, central, mutually orthogonal, and sum
    to the identity. Each is fixed by the conjugate-linear involution of
    $\C[G]$ sending $c\,g$ to $\bar c\,g^{-1}$.
\end{lemma}

\begin{proof}\leanok
    The first four statements hold for the identities of the factors of a
    product of matrix algebras. For the last, let $e=e_\lambda$, let $e^\ast$
    be its image under the involution, and let $P$ be the image of $e$ in the
    left regular representation, which is faithful and sends $e^\ast$ to the
    adjoint $P^\ast$. Then $e^\ast$ is central, so $e^\ast e$ is a central
    element of the factor $\lambda$ and equals $c\,e$ for a scalar $c$. Hence
    $P^\ast P=cP$. If $c=0$ then $P=0$, which is impossible. Taking adjoints
    gives $P^\ast=uP$ with $u=c/\bar c$, and idempotence of $P^\ast$ forces
    $u=1$. Thus $P^\ast=P$, and faithfulness gives $e^\ast=e$.
\end{proof}

\begin{definition}[Permutation operators and label projectors]
    \label{def:label_projectors}
    \lean{PermutationRepresentation.permOp,
        PermutationRepresentation.groupAlgebraRep,
        PermutationRepresentation.labelProj,
        PermutationRepresentation.labelObservable,
        PermutationRepresentation.labelEntropy}
    \leanok
    \uses{def:irrep_labels}
    Let $G$ act on a finite set $X$. The permutation operator $U(g)$ on $\C^X$
    sends the basis vector $e_x$ to $e_{gx}$, and extends linearly to an algebra
    homomorphism $\C[G]\to\mathrm{End}(\C^X)$. The label projector $\pi^\lambda$
    is the image of $e_\lambda$. For a real function $f$ of the labels, the
    label observable is $\sum_\lambda f(\lambda)\pi^\lambda$; for
    $f(\lambda)=\log d_\lambda$ it is the observable
    $F=\sum_\lambda(\log d_\lambda)\pi^\lambda$
    of~\cite[Section~6.1]{OpenAI2026AreaLaw}.
\end{definition}

\begin{lemma}[Label projectors are a resolution of the identity]
    \label{lem:label_projectors}
    \lean{PermutationRepresentation.labelProj_mul_labelProj,
        PermutationRepresentation.sum_labelProj,
        PermutationRepresentation.isHermitian_labelProj,
        PermutationRepresentation.commute_labelProj_of_forall_commute,
        PermutationRepresentation.isHermitian_labelObservable}
    \leanok
    \uses{def:label_projectors, lem:irrep_central_idempotents}
    The label projectors are mutually orthogonal orthogonal projections with sum
    the identity, and they commute with every operator that commutes with all
    permutation operators. Every label observable is self-adjoint.
\end{lemma}

\begin{proof}\leanok
    Apply the homomorphism to the relations of
    Lemma~\ref{lem:irrep_central_idempotents}. The image of the involution of
    $\C[G]$ is the adjoint, since $U(g)^\ast=U(g^{-1})$, so each projector is
    self-adjoint. An operator commuting with every $U(g)$ commutes with the
    image of every element of $\C[G]$.
\end{proof}

\section{Copies of a tensor power}

Let $\mathcal V=\bigotimes_{f\in F}\C^{n_f}$ be a finite-dimensional space
with a specified finite tensor factorization, and let $k\ge0$. For a subsystem
$Q\subseteq F$ and $\pi\in S_k$, the operator $U_Q(\pi)$ on
$\mathcal V^{\otimes k}$ permutes the $k$ copies of the factors in $Q$ and acts
as the identity on the other factors~\cite[Section~6]{OpenAI2026AreaLaw}.
Thus $U(\pi)=U_F(\pi)$ permutes entire copies, and the symmetric subspace
$\mathcal S_k=\Sym^k\mathcal V$ is the subspace fixed by every $U(\pi)$. The
Schur-label projectors $\pi^\lambda_{Q,k}$ are the label projectors of the
action $\pi\mapsto U_Q(\pi)$.

\begin{definition}[Copy permutations and the symmetric subspace]
    \label{def:copy_permutations}
    \lean{TensorPower.copyPerm, TensorPower.subsystemPerm,
        TensorPower.symmetricSubspace, PermutationRepresentation.symProj}
    \leanok
    \uses{def:label_projectors}
    The copy permutations of a subsystem, the symmetric subspace, and the
    averaging projector $\Pi_k=(k!)^{-1}\sum_{\pi\in S_k}U(\pi)$ onto it.
\end{definition}

\begin{lemma}[Central functions of $S_k$ are invariant under inversion]
    \label{lem:symmetric_central_inverse}
    \lean{IrrepLabel.coeff_conj_of_mem_center, IrrepLabel.coeff_inv_of_mem_center}
    \leanok
    The coefficients of a central element of $\C[G]$ form a class function.
    For $G=S_k$, a central element has the same coefficient at $\pi$ and at
    $\pi^{-1}$.
\end{lemma}

\begin{proof}\leanok
    Compare the coefficients of $h\,a$ and $a\,h$ at $hg$. A permutation and
    its inverse have the same cycle type, hence are conjugate.
\end{proof}

\section{The Schur-label estimates}

The following entries formalize the six parts of the Schur-label
lemma~\cite[Lemma 6.1]{OpenAI2026AreaLaw}.

\begin{theorem}[Schur-label estimates, part 2: commutation and complements]
    \label{thm:schur_label_commutation}
    \lean{TensorPower.commute_groupAlgebraRep_subsystemPerm_of_subset,
        TensorPower.commute_groupAlgebraRep_subsystemPerm_of_disjoint,
        TensorPower.commute_labelProj_subsystemPerm_of_subset,
        TensorPower.commute_labelProj_subsystemPerm_of_disjoint,
        TensorPower.groupAlgebraRep_mulVec_compl_of_mem_symmetricSubspace,
        TensorPower.labelProj_mulVec_compl_of_mem_symmetricSubspace,
        TensorPower.labelProj_mul_symProj_compl}
    \leanok
    \uses{def:copy_permutations, lem:label_projectors}
    Central label functions commute for nested or disjoint subsystems: if
    $Q\subseteq Q'$, a central function of the copy permutations of $Q$
    commutes with every operator generated by the copy permutations of $Q'$,
    and operators generated by the copy permutations of disjoint subsystems
    commute. On $\mathcal S_k$, complementary label projectors agree:
    $\pi^\lambda_{Q,k}\Pi_k=\pi^\lambda_{Q^c,k}\Pi_k$, even when $Q$ and $Q^c$
    have different dimensions.
\end{theorem}

\begin{proof}\leanok
    \uses{lem:symmetric_central_inverse}
    For $Q\subseteq Q'$ we have $U_{Q'}(\pi)=U_Q(\pi)U_{Q'\setminus Q}(\pi)$,
    where the second factor commutes with every operator generated by $Q$ and
    the first commutes with every central function of $Q$. Disjoint
    commutation is immediate. On $\mathcal S_k$,
    $U_Q(\pi)U_{Q^c}(\pi)z=z$, so $U_Q(\pi)z=U_{Q^c}(\pi^{-1})z$. Writing a
    central element as $\sum_\pi a(\pi)\pi$ with $a(\pi)=a(\pi^{-1})$ and
    reindexing by $\pi\mapsto\pi^{-1}$ shows that it acts in the same way
    through $Q$ and through $Q^c$ on $z$.
\end{proof}
